\documentclass[10pt,a4paper]{amsart}
\usepackage[margin=19mm]{geometry}
\usepackage{amsmath,amssymb,mathtools}
\usepackage[scr=rsfs,cal=euler]{mathalfa}
\usepackage{xfrac}
\usepackage[unicode,hidelinks,bookmarks=false]{hyperref}
\newcommand{\ie}{i.e.\ }
\newcommand{\cf}{cf.\ }
\newcommand{\resp}{respectively\ }
\newcommand{\Subsection}{\subsection}
\providecommand{\nobreakdash}{\nobreak\hbox{-}}
\setlength{\emergencystretch}{2em}
\multlinegap0pt
\usepackage{bm}
\usepackage{bbm}
\usepackage{dsfont}
\usepackage{xspace}
\usepackage{cancel}
\usepackage{mathtools}
\usepackage{amsthm}
\makeatletter
\@ifundefined{theorem}{\newtheorem{theorem}{Theorem}}{}
\@ifundefined{corollary}{\newtheorem{corollary}[theorem]{Corollary}}{}
\@ifundefined{lemma}{\newtheorem{lemma}[theorem]{Lemma}}{}
\makeatother
\title[Well-order a flame]{Well-order a flame}
\author{Zsuzsanna Jank\'o, Attila Jo\'o}
\date{Source: arXiv:2602.01184v1 (CC BY 4.0)}
\begin{document}
\maketitle

\noindent\emph{Source and licence.} Well-order a flame. Version arXiv:2602.01184v1 (CC BY 4.0), distributed under the Creative Commons Attribution 4.0 International licence, as recorded in the arXiv licence metadata for this exact version and shown on the abstract page https://arxiv.org/abs/2602.01184v1. Full rights notice: licenses/arxiv-2602.01184-CC-BY-4.0.txt.\par\smallskip

\noindent\textit{Editorial scope.} Counted unit: the Lemma whose source label is \texttt{lem: add fin} in the article (source0.tex lines 653--657, source label \texttt{lem: add fin}), delivered together with the article's own complete original proof of it (source0.tex lines 659--707, one \texttt{proof} environment of 49 lines). No mathematics is written, repaired or re-derived: statement and proof are the article's own text line for line. Required context retained uncounted: cor: at most n (lines 451--456); lem: char no single (lines 465--469); cor: insert one (lines 617--620). Every definition, display and label of those lines is retained unchanged. Omitted from this package and not redistributed: 1--450, 457--464, 470--616, 621--652, 658--658, 708--799. Nothing inside a retained excerpt is omitted or altered: every retained body is delivered exactly as the article's lines are written, so no omission receipt of any kind accompanies this package. Citations: the retained excerpts carry no citation command at all, so no bibliography entry of the article is transcribed and no reference is redistributed. Reference adaptation: none was needed and none was made; every reference inside the delivered unit resolves inside it, so no unresolved reference or citation is produced. Printed slips kept literal: no printed word, symbol, hypothesis or number of the counted statement or of its proof was altered. Layout and class: the article's own class is \texttt{amsart} with packages including amsmath, amssymb, amsthm, bm, mathtools, tikz, thmtools, float, xcolor, hyperref, biblatex, inputenc, csquotes, fontenc; the wrapper is the lane's pinned \texttt{amsart} editorial wrapper at 10pt on A4, which restates the theorem-like environments the retained text needs and adds the article's own macro definitions that the retained text needs (none), with extra wrapper packages (bm, bbm, dsfont, xspace, cancel, mathtools). The delivered numbering is the unit's own. Assets: no retained span contains a figure environment, a graphics-inclusion command, a PDF-page inclusion, a table or a TikZ picture; no image file is copied into this package, the package declares no asset dependency and no image file is redistributed with it. Licence: version arXiv:2602.01184v1 of this article is distributed under the Creative Commons Attribution 4.0 International licence, as recorded in the arXiv licence metadata for this exact version and shown on the abstract page https://arxiv.org/abs/2602.01184v1. A full rights notice is delivered at licenses/arxiv-2602.01184-CC-BY-4.0.txt. Characters: every retained excerpt is pure ASCII, so the delivered engine, XeLaTeX with its default Latin Modern text font, reports no missing character.
\begin{corollary}\label{cor: at most n}
Let $T$ be a $v$-tight set, and let $\mathcal{P}$ be a system of edge-disjoint $rv$-paths that covers all but $n\in\mathbb{N}$ edges of 
$\delta_D(v)$.
Let $\mathcal{P}'$ consist of the terminal segments of the paths in $\mathcal{P}$ starting from their last edges in $\delta_D(T)$.
Then at most $n$ edges of $\delta_D(T)$ are not used by the paths in $\mathcal{P}'$.
\end{corollary}

\begin{lemma}\label{lem: char no single}
Let $v\in V\setminus\{r\}$ and let $e\in\delta_D(v)$ be such that
$\delta_D(v)\setminus\{e\}\in\mathcal{G}_D(v)$.
Then $\delta_D(v)\notin\mathcal{G}_D(v)$ if and only if the tail $u$ of $e$ belongs to $T_{v,D-e}$.
\end{lemma}

\begin{corollary}\label{cor: insert one}
If $F\subseteq E$ is a flame, $v\in V\setminus\{r\}$, and $e\in E\setminus F$ such that
$e\in \delta_D(T_{v,F})$, then $F\cup\{e\}$ is a flame.
\end{corollary}

\begin{lemma}\label{lem: add fin}
If $F$ is an almost $\mathcal{P}$-respecting flame and $e\in E\setminus F$, then there
exists a finite set $\{e_i : i\leq n\}\subseteq E\setminus F$ containing $e$ such that
$F\cup\{e_i : i\leq m\}$ is a flame for every $m\leq n$.
\end{lemma}

\begin{proof}
Suppose, for a contradiction, that the statement is false, and let the quadruple
$(D,\mathcal{P},F,e)$ be a counterexample. Let $u$ and $v$ be the tail and head of $e$,
respectively. Define
\[
A:=\{e\}\cup\{f\in \delta_F(v) : P_f\not\subseteq F\}
\quad\text{and}\quad
A':=\bigcup_{f\in A} P_f .
\]
Since $F$ is almost $\mathcal{P}$-respecting,  $A$ is finite and thus so is  $A'$. We assume
that the counterexample is chosen so that
\[
n(D,\mathcal{P},F,e):=\lvert A'\setminus F\rvert
\]
is minimal.

To obtain a contradiction, it suffices to show that there exists an edge
$g\in A'\setminus F$ such that $F\cup\{g\}$ is a flame. Indeed, we must then have
$g\neq e$, since $(D,\mathcal{P},F,e)$ is a counterexample. Furthermore, replacing $F$ by
$F\cup\{g\}$ yields another counterexample with
\[
n(D,\mathcal{P},F\cup\{g\},e)=n(D,\mathcal{P},F,e)-1,
\]
contradicting the minimality of $n(D,\mathcal{P},F,e)$.

By Lemma~\ref{lem: char no single}, the set $T_{v,F}=:T$ contains $u$. By
Corollary~\ref{cor: insert one}, for any edge $g\in \delta_D(T)\setminus F$, the digraph
$F\cup\{g\}$ is a flame. Thus it remains to show that
\[
\delta_D(T)\cap(A'\setminus F)\neq\emptyset.
\]

Suppose, for a contradiction, that $\delta_D(T)\cap(A'\setminus F)=\emptyset$. Then for
every $f\in A$ we have $P_f\cap \delta_D(T)\subseteq F$. Note that the sets $P_f\cap \delta_D(T)=P_f\cap \delta_F(T)$ for $f\in A$
are nonempty and pairwise disjoint. Therefore,
\[
\left|\bigcup_{f\in A} P_f\cap \delta_F(T)\right|\geq |A|,
\]
and hence $|\delta_F(T)\cap A'|\geq |A|$.

Let
\[
\mathcal{Q}:=\{P_f : f\in \delta_F(v)\setminus(A\setminus\{e\})\}.
\]
On the one hand, $\mathcal{Q}$ covers all but $|A|-1$ ingoing edges of $v$ in $F$. On the
other hand, there are at least $|A|$ edges in
$\delta_F(T)$ that are not used by $\mathcal{Q}$ because $\bigcup_{f\in A} P_f\cap \delta_F(T)$ consists of such edges. This contradicts
Corollary~\ref{cor: at most n} applied to $F$  and the terminal segments of the paths in $\mathcal{Q}$ from their last edges in $\delta_{F}(T)$.
\end{proof}

\end{document}
